\section{Finite-level bounds for the ordinary branch parameters}
\label{sec:G88}

The branch degree in AC-G85 and AC-G86 can be bounded using the
original tame level. This removes an unspecified coefficient-field
degree from those estimates. It does not assert that the resulting
prime bound meets the original counting budget.

Throughout this section \(p=269\). Retain an actual ordinary cuspform
\(f_0\) supplied by AC-K29-good-twist and used in AC-G78: its weight
is \(2e+1\), its nebentypus is \(\eta_K\), and its tame level is
\(N\), with \(p\nmid N\). A chosen ordinary \(p\)-stabilization is
understood. The fixed quadratic character \(\tau\) in the
normalization of AC-G78 has values in \(\{1,-1\}\).
No assertion below requires the residual representation of \(f_0\)
to be irreducible.

Choose the unique integer \(w_0\) with
\[
 1\leq w_0\leq p-1,\qquad w_0\equiv e\pmod{p-1},
 \qquad k_0=2w_0-1.
 \tag{G88.1}
\]
Thus \(1\leq k_0\leq 2p-3\). Put
\[
 \begin{split}
 C(N,k_0)&=\dim_{\mathbf Q_p}
 H^1_{\mathrm{par}}\left(
 Y_1(Np)_{\overline{\mathbf Q}},
 \operatorname{TSym}^{k_0}\mathscr H_{\mathbf Q_p}(1)\right),\\
 \mu(Np)&=[\operatorname{SL}_2(\mathbf Z):\Gamma_1(Np)]
 =(Np)^2\prod_{\ell\mid Np}(1-\ell^{-2}).
 \end{split}
 \tag{G88.2}
\]
Here parabolic cohomology means the image of compactly supported
cohomology in the full \(H^1\). The notation is compatible with
the parabolic module in KLZ, Proposition 7.2.1.

\begin{theorem}[AC-G88: branch parameters bounded by the original level]
\label{thm:G88}
The ordinary branch used in AC-G78--AC-G86 may be taken over
\(B=\mathbf Z_p[[W]]\). Its domain \(R\) is a finite torsion-free
\(B\)-algebra; its normalization \(S\) is finite free over \(B\).
Write \(D=\operatorname{rank}_B S\). One can choose it with
\[
 D\leq \frac{C(N,k_0)}2
 \leq \mu(Np)(k_0+1)
 \leq 2p^3N^2.
 \tag{G88.3}
\]
The residue field of \(R\) or \(S\) need not be \(\mathbf F_p\).
Its degree is already included in \(D\); an extra coefficient-ring
extension is not needed for the trace quotient.

Consequently the quotient in AC-G85 can be chosen with
\[
 |Q_T|\leq p^{64D}\leq p^{32C(N,k_0)}
 \leq p^{128p^3N^2}.
 \tag{G88.4}
\]
The same bound applies after the delayed quadratic twist of AC-G81,
using the same branch and normalized coefficient algebra.

Retain all the actual occurrence, host, type and compatibility
hypotheses of AC-G83 and AC-G86. Let \(\Sigma,P,m\) have exactly
the meanings they have in AC-G86. Define
\[
 \mathcal N_N
   =48m(p-1)^2p^{\,128p^3N^2+4}.
 \tag{G88.5}
\]
Then the actual auxiliary prime may be chosen with
\[
 q\leq(\mathcal N_NP)^{310\mathcal N_N}.
 \tag{G88.6}
\]
One may use the sharper exponent \(32C(N,k_0)+4\) in (G88.5).
The primes in \(\Sigma\), including any fresh quadratic twist and
host primes, remain in \(P\) and \(m\). No bound for the original
Frame, EC, Big, counting-budget or whole-curve \(J\) conditions is
concluded.
\end{theorem}

\begin{proof}
\emph{The coefficient ring and the exact weight coordinate.}
Start with the full ordinary parabolic cohomology and its faithful
Hecke algebra over
\(\Lambda_D=\mathbf Z_p[[\mathbf Z_p^\times]]\), as in KLZ,
Proposition 7.2.1 and the paragraph immediately before Theorem
7.2.3. Since \(\mu_{p-1}\subset\mathbf Z_p^\times\) and \(p-1\)
is a unit, all its finite-character idempotents are defined over
\(\mathbf Z_p\). Select the component containing \(f_0\).

Let \(U=[1+p]-1\) be the literal source coordinate. The source
variable \(k\) is the coefficient exponent; a classical
specialization has modular weight \(k+2\).
Our half-weight variable \(w\) therefore satisfies \(k=2w-1\).
The exact change of coordinates is
\[
  U=(1+p)^{-1}(1+W)^2-1,
  \qquad W=(1+p)^w-1\ \text{at weight }w.
 \tag{G88.7}
\]
The constant term on the right lies in \(p\mathbf Z_p\), and its
linear coefficient is a unit. This gives an isomorphism
\(\mathbf Z_p[[U]]\simeq\mathbf Z_p[[W]]\) of complete local
rings. Explicitly its inverse is the square root congruent to
\(1\) of \((1+p)(1+U)\), minus \(1\); it exists uniquely because
\(2\) is a unit. If the source coordinate is first shifted from
\(k\) to \(k+1\), then \(U'=(1+p)(1+U)-1\) and the formula becomes
\(U'=(1+W)^2-1\).

Fixing \(w\equiv e\pmod{p-1}\), rather than only modulo
\((p-1)/2\), fixes the finite Teichmuller factor of the normalizing
cyclotomic character. For every \(a\in\mathbf Z_p^\times\), the
function
\[
 \omega(a)^{w_0}
 (1+W)^{\log\langle a\rangle/\log(1+p)}
 \in B^\times
 \tag{G88.8}
\]
is integral and specializes to \(a^w\) on that progression.
The binomial exponent belongs to \(\mathbf Z_p\).
Thus the Tate/half-weight normalization, the fixed quadratic
\(\tau\), and the character \(\eta_K\) all have their coefficient
values in \(B\). The continuous determinant and normalized trace
of AC-G78 are defined over the resulting \(B\)-Hecke branch.
The eigenvalues of \(f_0\) are encoded by its specialization of
that branch; they need not lie in \(\mathbf Q_p\).
There is no averaging over a tame diamond group whose order
might be divisible by \(p\). Instead use the whole tame Hecke
algebra before choosing the domain branch. A finite diamond
operator has finite order in that domain. If its image is
\(\alpha\), its value at \(f_0\) is the prescribed
\(\zeta=\eta_K(a)\in\{1,-1\}\). Factor
\(X^d-1=(X-\zeta)Q(X)\), where \(d\) is the operator's order.
The specialization of \(Q(\alpha)\) is \(d\zeta^{d-1}\ne0\)
in characteristic zero. Since the branch is a domain,
\((\alpha-\zeta)Q(\alpha)=0\) forces \(\alpha=\zeta\).
Thus its tame nebentypus is exactly \(\eta_K\) on the
whole branch, including when \(p\mid d\).
The cubic coefficient and the auxiliary finite local types may
require coefficient extensions later in the automorphic
construction. They are not adjoined to this rank-two trace
algebra.

\emph{Hecke rank from the actual ordinary module.}
Write \(\mathcal M\) for the selected finite-character component
of ordinary parabolic cohomology and \(\mathbb T\) for the
Hecke algebra acting faithfully on it. KLZ, Proposition
7.2.1(1), including its final parabolic clause, makes
\(\mathcal M\) finite projective over \(B\), hence free. Write
\(r=\operatorname{rank}_B\mathcal M\).
Already the faithful inclusion
\[
 \mathbb T\hookrightarrow\operatorname{End}_B(\mathcal M)
 \tag{G88.9}
\]
gives \(\operatorname{rank}_B\mathbb T\leq r^2\), without any
duality or Gorenstein assertion.

There is a stronger bound available from the same source.
The exact sequence in KLZ, Theorem 7.2.3 and its parts (ii),(v)
give, on this component,
\[
 \begin{gathered}
 0\longrightarrow F^+\mathcal M\longrightarrow\mathcal M
     \longrightarrow F^-\mathcal M\longrightarrow0,\\
 F^-\mathcal M\simeq\operatorname{Hom}_B(\mathbb T,B),
 \qquad F^+\mathcal M\simeq\mathbb T
       \quad\text{as }\mathbb T\text{-modules}.
 \end{gathered}
 \tag{G88.10}
\]
The second isomorphism need not be given a preferred basis.
The algebra \(\mathbb T\) is finite projective over \(B\), by the
paragraph preceding Theorem 7.2.3. Consequently
\[
 r=2\operatorname{rank}_B\mathbb T.
 \tag{G88.11}
\]
These statements precede the residual irreducibility and
distinction assumptions of KLZ, Theorem 7.2.8. No application of
that later theorem is made here.

Choose a minimal prime \(\mathfrak p\) of the local Hecke factor
below the characteristic-zero point \(f_0\), and let
\(R=\mathbb T/\mathfrak p\), using the corresponding factor of
\(\mathbb T\) if necessary. Minimal primes of a finite flat
algebra over the domain \(B\) contract to zero: a nonzero element
of \(B\) is regular on that algebra and belongs to none of its
minimal associated primes. Thus \(R\) is a finite torsion-free
domain over \(B\). Its fraction-field degree is at most
\(\operatorname{rank}_B\mathbb T\).
Normalization does not change this fraction field. The finite
normalization and freeness argument of AC-G84 therefore gives
\[
 D=\operatorname{rank}_B S
  =\operatorname{rank}_B R
  \leq\operatorname{rank}_B\mathbb T=r/2.
 \tag{G88.12}
\]
This argument permits nilpotents and repeated old roots in the
full Hecke algebra. It does not claim that algebra is reduced.

\emph{Perfect control at a bounded classical weight.}
In the literal source coordinate, KLZ, Proposition 7.2.1(4)
specializes at
\([1+p]-(1+p)^{k_0}\), with \(r=1\) in its depth notation.
Under (G88.7) this is the same specialization as
\(W=(1+p)^{w_0}-1\): the other factor of the difference of
squares is a unit. Perfect control, including the parabolic
clause and its finite-character action, identifies this
specialization of \(\mathcal M\) with an ordinary
finite-character summand of the cohomology in (G88.2).
Since \(\mathcal M\) is free, its specialized
\(\mathbf Q_p\)-dimension is \(r\). Therefore
\[
 r\leq C(N,k_0).
 \tag{G88.13}
\]
This uses a weight between \(3\) and \(2p-1\), irrespective of
the original size of \(e\). It does not assert that the
original eigenform itself has this smaller weight.

\emph{An elementary bound for the finite cohomology.}
Put \(M=Np\), so \(M\geq269\).
The group \(\Gamma_1(M)\) is torsion-free and does not contain
\(-I\). Indeed a finite-order matrix in
\(\operatorname{SL}_2(\mathbf Z)\) has trace in
\(\{-2,-1,0,1,2\}\), whereas its trace in \(\Gamma_1(M)\) is
\(2\pmod M\); only trace \(2\) is possible, and a finite-order
unipotent matrix is the identity.
Thus the complex curve is the quotient of the contractible
upper half-plane by this freely acting group.

For clarity, the coset index used in (G88.2) can be counted
directly. The stabilizer of the vector \((1,0)\) modulo \(M\)
is \(\Gamma_1(M)\). Every primitive vector modulo \(M\) has
a lift to a coprime pair of integers: after fixing a nonzero
representative of one coordinate, adjust the other by a
multiple of \(M\), using the Chinese remainder theorem at the
finitely many prime divisors of the first coordinate not
dividing \(M\). Bezout then completes the pair to a matrix
in \(\operatorname{SL}_2(\mathbf Z)\).
There are \(M^2\prod_{\ell\mid M}(1-\ell^{-2})\) primitive
vectors, by the Chinese remainder theorem. This proves the
index formula and the bound \(\mu(M)\leq M^2\).

The group \(\operatorname{SL}_2(\mathbf Z)\) is generated by
the usual matrices \(S,T\), by the Euclidean algorithm.
For a set of \(\mu(M)\) coset representatives, the Schreier
elements obtained by multiplying each representative by
\(S\) or \(T\) generate \(\Gamma_1(M)\); inverses are obtained
by reversing the same edges. Thus at most \(2\mu(M)\)
generators suffice. A crossed homomorphism to a
\((k_0+1)\)-dimensional coefficient space is determined by
its values on those generators. Its space, and therefore
its quotient \(H^1\), has dimension at most
\(2\mu(M)(k_0+1)\).

Apply the etale--complex comparison to the finite reductions
of the locally free coefficient sheaf, then pass to inverse
limits and tensor with \(\mathbf Q_p\).
The field extension from \(\overline{\mathbf Q}\) to
\(\mathbf C\) is legitimate for these torsion sheaves on a
curve. Precise references are Stacks, Tag 03SC(6), and Milne,
\emph{Lectures on Etale Cohomology} v2.21, Theorem 21.5,
with the inverse-limit convention in Section 19.
Parabolic cohomology is a subspace of \(H^1\), so
\[
 C(N,k_0)\leq 2\mu(Np)(k_0+1).
 \tag{G88.14}
\]
Combining (G88.12)--(G88.14), and using \(k_0+1\leq2p\),
proves (G88.3).

\emph{The trace quotient and delayed quadratic twist.}
The construction of AC-G85 now starts with
\(\mathcal O=\mathbf Z_p\). Its free order has rank \(4D\)
over \(B\), and reduction modulo \((p^4,W^4)\) has at most
\(p^{64D}\) elements. This proves (G88.4).
The possible residue extension in \(R\) or \(S\) has not
been discarded; it contributes to their \(B\)-rank \(D\).

Let \(\delta_0\) be the quadratic character used in AC-G81.
Use the same coefficient domain \(R\) and replace the
normalized determinant by its twist by \(\delta_0\), so its
trace is \(\delta_0T\) and its determinant is unchanged.
On a generic realization,
\(\rho_{\delta_0}(g)=\delta_0(g)\rho(g)\).
The two \(S\)-linear spans of the group matrices are identical,
because each scalar \(\delta_0(g)\) is \(1\) or \(-1\).
Their reflexive hull over \(B\), used in AC-G85, is therefore
the same free order. The quotient cardinality bound is
unchanged. The representation may acquire ramification at
the fresh prime \(r_0\), which must be included in \(\Sigma\).
The dense classical family consists of the actual quadratic
twists of the original family; its weight and nebentypus are
unchanged. This does not identify it with an old Dat after a
parity-changing twist, nor require using a larger Hecke
branch to bound its coefficients.

Finally apply AC-G86 with \(c=1\) and \(D\leq2p^3N^2\).
Its degree bound is
\(48m(p-1)^2p^{64D+4}\), bounded by (G88.5).
Monotonicity of \(x\mapsto(xP)^{310x}\) for \(x\geq1\)
gives (G88.6). The normalized-\(S\) precision of AC-G84 is
exactly the precision delivered by AC-G85; no equality
\((p^4,W^4)S\cap R=(p^4,W^4)R\) is used.
All later actual Hecke-specialization and local-type
hypotheses remain those already checked in AC-G78--AC-G86.
\end{proof}

\paragraph{Reference scope and geometric dictionary.}
The source input for ordinary control and the two ordinary
pieces is KLZ v5, Section 7.1, Proposition 7.2.1 including its
parabolic clause, the following finite-projectivity paragraph,
and Theorem 7.2.3(ii),(v) with its complete printed normalization
proof, printed pp.~35--36. The imported Ohta theorems are
identified there; they have not been re-proved in this section.
The source comparisons are Stacks Tag 03SC(6), with its printed
proof, and Milne v2.21, Section 21, especially Theorem 21.5
and its printed proof through Lemma 21.6, together with the
finite-cover comparison in Theorem 21.3.
They identify the finite-level local systems used to estimate
cohomology. The passage from a finite Hecke algebra to a
minimal domain and its normalization is a finite morphism of
weight spaces. Only this coefficient-size calculation is
being made; no density assertion about a prescribed
deformation curve, and no new Grothendieck theorem, is
attributed to it.
